\section{Combat Results}

\begin{rulebox}{General Rule}
There are six possible combat results: Ad (Attacker Disrupted), Dd (Defender Disrupted), Dx (Disruption Exchange), De (Defender Eliminated), Ae (Attacker Eliminated), and a dash (No Effect). Only one of these results will occur for each combat. Results are applied to the affected units and leaders immediately, before any other combats are resolved.
\end{rulebox}

\ruleslabel{Procedure}

The active player looks at the Combat Table and reads across the row of entries indicated by the combat-resolution die roll until he comes to the column for the combat ratio. Where row and column meet, he finds a result whose effects are applied immediately to the affected units.

\ruleslabel{Cases}

\caseheading{11.1} If the result is an Ad (Attacker Disrupted), all attacking units are immediately disrupted. Disrupted units are governed by special movement and combat considerations detailed in cases 11.7 and 11.8. Disrupted units remain disrupted until rallied as detailed in 6.0 (How to Rally).

\caseheading{11.2} If the result is a Dd (Defender Disrupted), all defending ordered units in the combat are immediately disrupted.

Disrupted defending units in the combat are immediately eliminated (removed from play) as if they had suffered a De result (11.4). Otherwise, disrupted units are governed by special movement and combat rules (11.7 and 11.8). Disrupted units remain disrupted until rallied as detailed in 6.0 (How to Rally).

\caseheading{11.3} If the result is a Dx (Disruption Exchange), both the attacker and the defender suffer disruption.

First, all disrupted defending units are eliminated (as if they had suffered the De result described in 11.4). Second, all ordered defending units are disrupted. Finally, the active player disrupts a number of his own units whose printed combat strength at least equals the face-up combat strength of the disrupted or eliminated defending units. Only units actually participating in the combat are affected by these results. If the total printed combat strengths of the attacking units are less than the face-up combat strengths of the defending units, all attacking units are simply disrupted; there is no further effect. Disrupted units can be rallied as detailed in 6.0 (How to Rally).

\caseheading{11.4} If the result is a De (Defender Eliminated), all the defending units are immediately and permanently eliminated (removed from play). Eliminated units count for demoralization and victory-point purposes as detailed in 13.0 (Demoralization) and 15.0 (How to Win). They have no further effect on play and never re-enter the game.

\caseheading{11.5} If the result is an Ae (Attacker Eliminated), all the attacking units are immediately and permanently eliminated from play. Eliminated units count for demoralization and victory-point purposes as detailed in 13.0 (Demoralization) and 15.0 (How to Win). They have no further effect on play and never re-enter the game.

\caseheading{11.6} If the result is a dash (No Effect), there is no effect whatsoever from the combat, and the players proceed to resolve the next combat situation.

\caseheading{11.7} Disrupted units are flipped to their back sides. Disrupted units have their combat strengths and movement allowances reduced to the numbers printed on their back sides until they are rallied.

\caseheading{11.8} Disrupted units suffer special movement and combat restrictions:

\begin{tightitems}
  \item Disrupted units cannot attack. They defend normally, but if they suffer a second disruption effect as a result of combat (not artillery fire) while disrupted, they are eliminated.
  \item Disrupted units have no ZOCs, and the active player is never required to attack them (though he can do so at his option).
  \item Disrupted foot units can be charged by Heavy Horse units.
  \item Disrupted units are unaffected by enemy ZOCs.
  \item Disrupted units do not expend MPs to move. Instead, the movement allowance on the back of each unit is the number of hexes the unit can move each friendly movement phase while disrupted. Disrupted units can enter all types of terrain and cross all types of hexsides that they would be permitted to enter or cross while ordered, regardless of the normal MP cost (see also 13.7).
\end{tightitems}

Disruption results only affect units. Leaders and artillery are not affected in any way by the disruption of a unit with which they are stacked.

\caseheading{11.9} Leaders are not directly affected by combat results. However, if the unit with which a leader is stacked is eliminated, the player who controls that leader must make an evasion roll to survive.

The player counts the distance in hexes from the leader to the nearest friendly unit, ignoring the presence of intervening terrain or enemy units. Distances greater than 6 are treated as 6. The player then rolls a single six-sided die. If the result is greater than or equal to the distance to the nearest friendly unit, the leader is placed atop that unit. If there are two or more friendly units onto which a leader could displace, the player controlling the leader decides to which hex the leader goes. If the die roll is less than this distance, the leader is immediately and permanently eliminated (removed from play).

Leaders thus eliminated count for 13.0 (Demoralization) and 15.0 (Victory Point) purposes. There is no limit to the number of times a leader can be forced to displace if stacked with units that are subsequently eliminated; however, the leader does not add his command rating to any unresolved combat in a phase in which the leader has already been displaced.
